\scanpage{77}
% archive foliation: 76
On a aussi : $B'(f,g)=\dfrac{dg}{g}-\dfrac{df}{f}$, \unsure{\ldots}

Si on veut par contre un invariant sur $E_X^{\deleted{*}}$, on prendra
\[
g\longrightarrow\frac{g}{dg}\qquad
E_X\longrightarrow\underline\omega_{X/\mathbf C}^{\otimes(-1)}
\]
on \deleted{trouve} \unsure{définit}
\[
\unsure{C|B}^{\unsure{!|\prime}}(f,g)=
\deleted{\frac{dg}{g}-\frac{df}{f}}
g(dg)^{-1}-f(df)^{-1}
\]
satisfaisant à a) b) c)
$[C^{\unsure{?|\prime}}(f,g)=C^{\unsure{?|\prime}}(f,ag)$
si $a\in\mathbf C^*)$
d) $(C^{\unsure{?|\prime}}(f,g)=C^{\unsure{?|\prime}}(f,h)$ et
$g(x)=h(x)$ en \uline{un pt qui n'est pas un zéro} de $g$ ou $h$
$\Longrightarrow g=h]$ et e) \deleted{$g\longmapsto C(f,g)$}
\unsure{pas tenu|pas très} : \unsure{faux e|fait e}, car
\deleted{$g(dg)^{-1}-f(df)^{-1}$}
p.ex. $g(dg)^{-1}$ ne peut être id. nul si $g\in\Gamma E_X$ !
Donc il faut se borner à
$E_X^*\inserted{=\widetilde{\operatorname{Et}}_{\mathbf C.\mathrm{an}}
(X,\mathbf C^*)}$, auquel cas on l'envoie dans
$\underline\omega_{X/\mathbf C}^{\otimes(-1)*}$ par $g(dg)^{-1}$,
ou aussi dans $\underline\omega_{X/\mathbf C}^*$ par $dg/g$,
et il y a lieu de considérer $\omega_{X/\mathbf C}^*$ comme un torseur
sous $\underline{\mathcal O}_X^*$.

Donc ici on prendra
\[
C^{\unsure{?|\prime}}(f,g)=\frac{dg}{g}:\frac{df}{f}
=\frac{dg}{df}\frac{g}{f}\in\Gamma(X,\underline{\mathcal O}_X^*)
\]

Dans tous ces cas, l'espace des valeurs \unsure{par|prises par} les $f,g$
\inserted{$(\mathbf P^1(\mathbf C),\underline E^1(\mathbf C),G_m(\mathbf C))$}
\unsure{étant} un \uline{espace homogène} sous le groupe d'op.
envisagé $(\unsure{GP(1)|Gb(1)},G_a,G_m)$ ce qui semble
\unsure{condition} de succès ($G_m$ opérant sur $\mathbf C$ : mauvais !)

c) Groupe $\operatorname{Aff}(1)$ \deleted{$\unsure{G}$}
$=G_m\mathbin{.}G_a$ opère sur $\underline E^1$.

On considère donc
\figure{\begin{tikzpicture}[x=1cm,y=1cm]
\draw (0,0)--(0,0.58);\draw (0.25,0.04)--(0.24,0.61);
\end{tikzpicture}}

